% !PRINTED 7
\weberpage{7}
%<<CONT
then the identity holds:
\weq{4}{a+b+c=0}
and we obtain the six cross ratios
\weq{5}{\begin{aligned}
&-\frac{c}{b},\quad -\frac{a}{c},\quad -\frac{b}{a},\\
&-\frac{b}{c},\quad -\frac{c}{a},\quad -\frac{a}{b};
\end{aligned}}
if we set the first of them $-c:b=\varkappa^2$, then by (4) one obtains the six:
\weq{6}{\varkappa^2,\quad 1-\frac{1}{\varkappa^2},\quad \frac{1}{1-\varkappa^2},\quad \frac{1}{\varkappa^2},\quad \frac{\varkappa^2}{\varkappa^2-1},\quad 1-\varkappa^2.}

Thus the six cross ratios formed from the four points are expressed \gesperrt{linearly} by one among them.

If the $x_1, x_2, x_3, x_4$ are permuted among themselves, then the $a, b, c$ are permuted among themselves and change their signs, yet in such a way that either all three signs are changed at once or all remain unchanged. For example the transpositions

\[
\begin{aligned}
(14)\ \text{and}\ (23) & \begin{pmatrix} a & b & c \\ -a & -c & -b \end{pmatrix} \\
(24)\ \text{and}\ (31) & \begin{pmatrix} a & b & c \\ -c & -b & -a \end{pmatrix} \\
(34)\ \text{and}\ (12) & \begin{pmatrix} a & b & c \\ -b & -a & -c \end{pmatrix}
\end{aligned}
\]

% [sic] The comma after x_4 is printed; expected no comma before "are real".

If $x_1, x_2, x_3, x_4$, are real, then $\varkappa^2$ is a positive proper fraction then and only then, when $\varkappa^2$ and $1-\varkappa^2$ are positive, that is, when either $a$ and $c$ are positive and $b$ negative, or $a$ and $c$ negative and $b$ positive. This occurs when $x_1, x_2, x_3, x_4$ follow one another in this order, ascending in magnitude, and in all those arrangements which arise from it by such permutations as leave $b$ unchanged. These are the following eight:
\[
\begin{array}{cccc}
1 & 2 & 3 & 4 \\
2 & 3 & 4 & 1 \\
3 & 4 & 1 & 2 \\
4 & 1 & 2 & 3 \\
1 & 4 & 3 & 2 \\
4 & 3 & 2 & 1 \\
3 & 2 & 1 & 4 \\
2 & 1 & 4 & 3
\end{array}
\]
